% ── Prototipado ──────────────────────────────────────────────────────────────
\section{Prototipado}

\subsection{Prototipado Dinámico}

\subsubsection{Descripción y propósito}

El prototipado dinámico constituye una representación interactiva del sistema a
desarrollar, elaborada antes de la implementación definitiva del código. Su propósito
principal es visualizar el flujo de navegación, la distribución de los elementos de
la interfaz y la experiencia general del usuario, permitiendo detectar inconsistencias
de diseño y validar los requerimientos funcionales en una etapa temprana del proyecto.

A diferencia de los mockups estáticos, el prototipo dinámico permite simular la
interacción del usuario con la plataforma ---clics en botones, transiciones entre
vistas, formularios de registro y flujos de consulta--- sin necesidad de código
funcional. Esto facilita la retroalimentación del docente y de los usuarios potenciales
antes de comprometer esfuerzo de desarrollo, reduciendo el riesgo de construir
interfaces que no satisfagan las necesidades reales de los actores del sistema.

\subsubsection{Herramienta utilizada}

El prototipo dinámico se elaborará utilizando \textbf{Figma}, herramienta estándar del
sector para diseño de interfaces y prototipado interactivo. Figma permite crear
componentes reutilizables, establecer enlaces de navegación entre pantallas y compartir
el prototipo mediante un enlace web para su revisión, sin instalación de software
adicional por parte del evaluador.

\subsubsection{Vistas y flujos prototipados}

El prototipo dinámico contempla las siguientes pantallas y flujos de interacción
principales:

\begin{table}[H]
\centering
\small
\caption{Vistas contempladas en el prototipo dinámico}
\begin{tabular}{C{0.5cm} L{3.8cm} L{8.8cm}}
\toprule
\rowcolor{unsaacRed}
\color{white}\textbf{N.º} &
\color{white}\textbf{Vista / Módulo} &
\color{white}\textbf{Descripción} \\
\midrule
1 & Página de inicio &
Portal principal de la plataforma con acceso a la consulta del TUPA,
inicio de sesión y novedades institucionales. \\
\addlinespace[3pt]
2 & Consulta del TUPA &
Buscador y listado de procedimientos con filtros por tipo, área
responsable y costo; detalle de cada procedimiento con requisitos y
plazos. \\
\addlinespace[3pt]
3 & Registro / Inicio de sesión &
Formulario de acceso para estudiantes y personal administrativo, con
opción de recuperación de contraseña. \\
\addlinespace[3pt]
4 & Registro de solicitud &
Formulario paso a paso para iniciar un trámite, incluyendo la carga
de documentos digitales adjuntos. \\
\addlinespace[3pt]
5 & Seguimiento de trámites &
Panel del usuario con historial de solicitudes, estado actual
(recibido, en proceso, observado, resuelto) y línea de tiempo del
trámite. \\
\addlinespace[3pt]
6 & Centro de notificaciones &
Vista de alertas sobre cambios de estado en los trámites del usuario,
con marcadores de leído/no leído. \\
\addlinespace[3pt]
7 & Panel administrativo &
Bandeja de trámites recibidos para el personal de oficina, con
herramientas de validación documental y cambio de estado. \\
\addlinespace[3pt]
8 & Módulo de reportes &
Gráficas y tablas con estadísticas de trámites procesados, filtrables
por período, tipo y estado de resolución. \\
\bottomrule
\end{tabular}
\end{table}

\subsubsection{Flujos de interacción prototipados}

El prototipo dinámico representará los dos flujos principales del sistema:

\medskip
\textbf{Flujo del usuario estudiante:}

\begin{center}
\begin{tikzpicture}[
  node distance=1.0cm,
  box/.style={
    draw=unsaacRed, fill=unsaacRed!8, rounded corners=3pt,
    text width=2.3cm, align=center, font=\tiny, minimum height=0.7cm
  },
  arr/.style={->, >=stealth, draw=unsaacGray, thick}
]
  \node[box] (A) {Inicio de sesión};
  \node[box, right=of A] (B) {Búsqueda de procedimiento};
  \node[box, right=of B] (C) {Consulta de requisitos};
  \node[box, right=of C] (D) {Registro de solicitud};
  \node[box, right=of D] (E) {Seguimiento del trámite};
  \draw[arr] (A)--(B);
  \draw[arr] (B)--(C);
  \draw[arr] (C)--(D);
  \draw[arr] (D)--(E);
\end{tikzpicture}
\end{center}

\medskip
\textbf{Flujo del administrador:}

\begin{center}
\begin{tikzpicture}[
  node distance=1.0cm,
  box/.style={
    draw=unsaacGray, fill=lightGray, rounded corners=3pt,
    text width=2.4cm, align=center, font=\tiny, minimum height=0.7cm
  },
  arr/.style={->, >=stealth, draw=unsaacGray, thick}
]
  \node[box] (A) {Acceso al panel};
  \node[box, right=of A] (B) {Revisión de solicitudes};
  \node[box, right=of B] (C) {Validación de documentos};
  \node[box, right=of C] (D) {Cambio de estado};
  \node[box, right=of D] (E) {Cierre del trámite};
  \draw[arr] (A)--(B);
  \draw[arr] (B)--(C);
  \draw[arr] (C)--(D);
  \draw[arr] (D)--(E);
\end{tikzpicture}
\end{center}

\subsubsection{Alcance dentro del primer entregable}

En el marco del presente primer entregable, el prototipo dinámico se presenta como
una propuesta inicial de diseño que acompaña al plan de proyecto. Su alcance incluye:

\begin{itemize}
  \item La definición de la paleta de colores, tipografía y guía de estilo visual de
        la plataforma, con identidad afín a los colores institucionales de la UNSAAC.
  \item El diseño de las vistas descritas en la tabla anterior, con navegación funcional
        entre pantallas mediante los prototipos interactivos de Figma.
  \item La representación del flujo principal del usuario estudiante y del flujo del
        administrador, tal como se detallan en los diagramas anteriores.
\end{itemize}

El prototipo no constituye la implementación final de la plataforma; su función es
servir como artefacto de comunicación y validación entre el equipo de desarrollo, el
docente y los usuarios potenciales, sentando las bases para el diseño definitivo que
se elaborará en el segundo entregable semestral.

\subsubsection{Enlace al prototipo}

\begin{supuesto}
\small\textbf{Propuesta inicial:} El enlace al prototipo interactivo en Figma será
incluido aquí una vez concluido su diseño. Se recomienda publicarlo en modalidad
\emph{View Only} para facilitar la revisión por parte del docente y los usuarios
de prueba:\\[4pt]
\texttt{[Enlace al prototipo Figma --- por completar antes de la entrega formal]}
\end{supuesto}
